\documentclass{article}
\usepackage{ctex}

\usepackage{amsmath}    % 数学公式排版
\usepackage{amssymb}    % 数学符号
\usepackage{amsthm}     % 定理证明格式

\usepackage{geometry}   % 设置页边距
\geometry{a4paper, margin=2.54cm}

\usepackage{setspace}   % 设置行距
\onehalfspacing         % 1.5 倍行距

\usepackage[colorlinks]{hyperref}


\begin{document}

\section{带编号的行间公式及其引用示例}

在量子力学中，微观粒子在势场 \(V(\mathbf{r})\) 中的运动可以由方程\eqref{eq:schrodinger-equation}描述：
\begin{equation}
    \label{eq:schrodinger-equation}
    \mathrm{i} \hbar \frac{\partial \Psi}{\partial t} = 
    - \frac{\hbar^2}{2m} \nabla^2 \Psi + V(\mathbf{r}) \Psi
\end{equation}
其中，$m$是粒子质量，$\hbar$是约化普朗克常数，$\Psi=\Psi(\mathbf{r},t)$是粒子的波函数。
方程\eqref{eq:schrodinger-equation}被称为\textbf{薛定谔方程}。


\section{多行公式排版}

\subsection{多行公式用同一编号}

Maxwell电磁方程组如\eqref{eq:maxwell-equations}所示。
\begin{equation}
    \label{eq:maxwell-equations}
    \left\{
      \begin{aligned}
        \nabla \cdot \mathbf{D} &= \rho_0 \\
        \nabla \cdot \mathbf{B} &= 0 \\
        \nabla \times \mathbf{E} &= -\frac{\partial \mathbf{B}}{\partial t} \\
        \nabla \times \mathbf{H} &= \mathbf{j}_0 + \frac{\partial \mathbf{D}}{\partial t}
      \end{aligned}
    \right.
\end{equation}


\subsection{多行公式每行单独编号}
Maxwell电磁方程组如\eqref{eq:Maxwell-Eqns-1}--\eqref{eq:Maxwell-Eqns-4}所示。
\begin{align}
    \label{eq:Maxwell-Eqns-1}
    \nabla \cdot \mathbf{D} &= \rho_0 \\
    \label{eq:Maxwell-Eqns-2}
    \nabla \cdot \mathbf{B} &= 0 \\
    \label{eq:Maxwell-Eqns-3}
    \nabla \times \mathbf{E} &= -\frac{\partial \mathbf{B}}{\partial t} \\
    \label{eq:Maxwell-Eqns-4}
    \nabla \times \mathbf{H} &= \mathbf{j}_0 + \frac{\partial \mathbf{D}}{\partial t}
\end{align}

我们也可以用 \verb|subequations| 环境使不同方程拥有公共编号和各自的子编号。例如：
\begin{subequations}
    \begin{align}
        \label{eq:Maxwell-Eqns-01}
        \nabla \cdot \mathbf{D} &= \rho_0 \\
        \label{eq:Maxwell-Eqns-02}
        \nabla \cdot \mathbf{B} &= 0 \\
        \label{eq:Maxwell-Eqns-03}
        \nabla \times \mathbf{E} &= -\frac{\partial \mathbf{B}}{\partial t} \\
        \label{eq:Maxwell-Eqns-04}
        \nabla \times \mathbf{H} &= \mathbf{j}_0 + \frac{\partial \mathbf{D}}{\partial t}
    \end{align}
\end{subequations}


\subsection{分段函数示例}

可使用 \verb|cases| 环境产生分段函数，例如：
\begin{equation}
    f(x,y)=
    \begin{cases}
    \dfrac{xy}{x^2+y^2}, & x^2+y^2 \neq 0 \\
    0, & x=y=0
    \end{cases}
\end{equation}
% 思考：将 \dfrac 命令改为 \frac，排版效果会有怎样的不同？
%       如果放在正文中（行内公式）中，两者有何不同？


    


\end{document}